\section{System Model and Problem Formulation}

\subsection{Multi-tier LMC System Model}

We consider a multi-tier LMC system consisting of a set of satellites $\mathcal{S}$, local controllers at ground stations, and a centralized LMC controller (LMCC) operating on a cloud server~\cite{hu2022software}. The LMC network is organized into multiple tiers $\mathcal{K}$, where each tier $k \in \mathcal{K}$ comprises orbits at the same altitude. Each satellite $i \in \mathcal{S}$ in tier $k$ moves along its respective orbital plane. Let $\mathcal{N}_i$ denote the set of satellites directly connected to satellite $i$ via intra-tier or inter-tier ISLs, where $j \in \mathcal{N}_i$ if the angular separation $\phi_{ij} \le \phi_{\mathrm{beam}}$ and $\phi_{\mathrm{beam}}$ is the beamwidth of the satellite antenna. The neighbor set $\mathcal{N}_i$ is further categorized into an intra-tier set $\mathcal{N}_i^{\mathrm{intra}}$ and an inter-tier set $\mathcal{N}_i^{\mathrm{inter}}$ based on their respective orbital altitudes. Let $\mathcal{R}$ denote the set of source--destination pairs served by the LMC network.
The LMCC maintains time-varying network snapshots generated from the two-line element (TLE) datasets, which provide satellite altitude $h_i$, inclination $\alpha_i$, and phase offset $\mu_i$. We assume that the network topology remains static within each snapshot interval $T$~\cite{han2023time}. \textcolor{blue}{However, ISL states including $l_{ij}$ and $c_{ij}$ are recomputed from the current satellite positions and velocities at each snapshot.}

As shown in Fig.~\ref{fig:system}, the procedure of the proposed algorithm at the beginning of each snapshot interval is as follows:
(1) Each satellite $i$ measures the signal-to-interference-plus-noise ratio (SINR) for its ISLs and transmits these measurements to the LMCC via its connected local controller.
(2) The LMCC calculates the duration ($l_{ij}$) and capacity ($c_{ij}$) of each ISL and computes the ISL utility ($U_{ij}$).
(3) The LMCC performs GNN-based link pruning on $\mathcal{G}$ to construct a pruned graph $\mathcal{G}'$.
(4) The LMCC determines routing paths ($P_{sd}^*$) on $\mathcal{G}'$ by a utility-based heuristic to maximize the \textcolor{blue}{average ISL utility}.
(5) Finally, the LMCC distributes the routing paths ($P_{sd}^*$) to the respective satellites through the local controllers.
\textcolor{blue}{The communication overhead is bounded by $|\mathcal{S}||\mathcal{N}_i|$ SINR reports per snapshot interval $T$, and each satellite receives only its own routing path $P_{sd}^*$ after path determination, so both uplink and downlink control overhead are negligible compared to $T$~\cite{han2023time},~\cite{hu2022software}.}

\begin{figure}[t]
\centering
\includegraphics[width=1\linewidth]{Image/SystemModel.PNG}
\caption{Overall procedure of the proposed algorithm.}
\label{fig:system}
\end{figure}

\subsubsection{ISL Duration}

Let $l_{ij}$ denote the ISL duration, defined as the time interval during which a link between satellites $i$ and $j$ remains available for data transmission~\cite{han2023time}. The position vector $\mathbf{p}_i$ of satellite $i$ is given by
\begin{equation} \label{eq:ecef_pos}
\mathbf{p}_i = (R + h_i) \left[\cos \varphi_i \cos \lambda_i,~ \cos \varphi_i \sin \lambda_i,~ \sin \varphi_i \right]^\top,
\end{equation}
where $R$ is the Earth's radius, and $\varphi_i$ and $\lambda_i$ are the geocentric latitude and longitude of satellite $i$, respectively. The ISL distance between satellites $i$ and $j$ is $d_{ij} = \|\mathbf{p}_j - \mathbf{p}_i\|$.

Let $\mathbf{v}_i$ and $\mathbf{v}_j$ denote the velocity vectors of satellites $i$ and $j$, respectively. The rate of change of the ISL distance $\dot{d}_{ij}$ is given by
\begin{equation} \label{eq:dist_rate}
\dot{d}_{ij} = \frac{\mathbf{v}_{ij} \cdot \mathbf{r}_{ij}}{\|\mathbf{r}_{ij}\|} = \|\mathbf{v}_{ij}\| \cos(\theta_{ij}),
\end{equation}
where $\mathbf{v}_{ij} = \mathbf{v}_j - \mathbf{v}_i$ is the relative velocity vector, $\mathbf{r}_{ij} = \mathbf{p}_j - \mathbf{p}_i$ is the relative position vector, and $\theta_{ij}$ is the angle between $\mathbf{v}_{ij}$ and $\mathbf{r}_{ij}$. Assuming a constant relative velocity during a short interval, the ISL duration $l_{ij}$ is given by
\begin{equation} \label{eq:link_duration}
l_{ij} = \begin{cases}
    \left| \frac{d_{ij}}{\dot{d}_{ij}} \right|, & \text{if } |\varphi_i|, |\varphi_j| \leq \varphi_p, \\
    0, & \text{otherwise},
\end{cases}
\end{equation}
where $\varphi_p$ is the polar latitude threshold.
\textcolor{blue}{Here, $|d_{ij}/\dot{d}_{ij}|$ represents the time until the link between satellites $i$ and $j$ becomes unavailable due to relative motion.}

\subsubsection{ISL Capacity}

The angular separation $\phi_{ij}$ is defined as the angle between the nadir vector of satellite $i$ and the ISL direction toward satellite $j$. Using the law of cosines, $\phi_{ij}$ is given by
\begin{equation} \label{eq:angular_separation}
\phi_{ij} = \arccos \left[\frac{(R+h_i)^2 + d_{ij}^2 - (R+h_j)^2}{2(R+h_i)d_{ij}}\right].
\end{equation}
When satellite $i$ transmits data to satellite $j$, the receiver $j$ experiences interference from other neighboring satellites~\cite{cao2025dense}. The intra-tier interference $I_j^{\mathrm{intra}}$ at satellite $j$ is given by
\begin{equation} \label{eq:interference_intra}
I_j^{\mathrm{intra}} = \sum_{j^{\prime} \in \mathcal{N}_{j}^{\mathrm{intra}} \setminus \{i\}} \frac{P_o G_t G_r b_{j^{\prime}j}^2}{PL(d_{j^{\prime}j})},
\end{equation}
where $P_o$ is the transmit power, $G_t$ and $G_r$ are the transmitter and receiver antenna gains, $PL(d_{j^{\prime}j})$ is the free-space path loss, and $b_{j^{\prime}j}$ is the beam misalignment factor. The inter-tier interference $I_j^{\mathrm{inter}}$ is given by
\begin{equation} \label{eq:interference_inter}
I_j^{\mathrm{inter}} = \sum_{j^{\prime} \in \mathcal{N}_j^{\mathrm{inter}}\setminus \{i\}} \eta(h_j,h_{j^{\prime}}) \frac{P_o G_t G_r b_{j^{\prime}j}^2}{PL(d_{j^{\prime}j})},
\end{equation}
where $\eta(h_j,h_{j^{\prime}}) = \exp(-\kappa |h_j-h_{j^{\prime}}|)$ is the altitude-dependent attenuation factor and $\kappa$ is the attenuation coefficient. \textcolor{blue}{The factor $\eta$ is applied only to inter-tier interferers, which arrive from off-boresight directions where the antenna gain is reduced~\cite{cao2025dense}.} The SINR for the ISL from satellite $i$ to $j$ is given by~\cite{cao2025dense}
\begin{equation} \label{eq:sinr}
\mathrm{SINR}_{ij} = \frac{P_o G_t G_r b_{ij}^2}{PL(d_{ij})\left(\textcolor{blue}{N_0 B} + I_{j}^{\mathrm{intra}} + I_{j}^{\mathrm{inter}}\right)},
\end{equation}
where $N_0$ is the noise power spectral density\textcolor{blue}{, $B$ is the available bandwidth, and $N_0 B$ denotes the total noise power}. The ISL capacity $c_{ij}$ is given by
\begin{equation}\label{eq:link_capacity}
    c_{ij} = B \log_2(1 + \mathrm{SINR}_{ij}).
\end{equation}

\subsection{Link Utility Model and Problem Formulation}

We model the LMC network as a graph $\mathcal{G} = (\mathcal{S}, \mathcal{E})$, where $\mathcal{E}$ denotes the set of ISLs. The quality of each ISL $(i,j) \in \mathcal{E}$ is quantified by a duration utility $U_{ij}^l$ normalized by the snapshot interval $T$ and a capacity utility $U_{ij}^c$ normalized via a softmax function\textcolor{blue}{, which suppresses the outlier sensitivity of min-max normalization and provides a stable relative ranking of link capabilities across all ISLs~\cite{fu2024softact}}, given by
\begin{equation} \label{eq:utility_components}
    U_{ij}^l = \frac{l_{ij}}{T}, \quad U_{ij}^c = \frac{\exp(c_{ij})}{\sum_{(m,n) \in \mathcal{E}} \exp(c_{mn})}.
\end{equation}
The ISL utility $U_{ij}$ is then defined as
\begin{equation} \label{eq:utility_def}
U_{ij} = \beta \cdot U_{ij}^l + (1-\beta) \cdot U_{ij}^c,
\end{equation}
where $\beta \in [0,1]$ is a weighting factor. \textcolor{blue}{Since $l_{ij}$ represents the time until the link becomes unavailable~\cite{han2023time} and $c_{ij}$ is determined by $\mathrm{SINR}_{ij}$ via Shannon's formula in Eq.~\eqref{eq:link_capacity}, ISLs with higher $U_{ij}$ are less likely to fail and provide higher achievable rate, both of which directly improve throughput~\cite{wang2023reliability}.}

Let $x_{ij}^{sd} \in \{0, 1\}$ denote a binary variable indicating whether ISL $(i,j) \in \mathcal{E}$ is included in the routing path for source--destination pair $(s,d) \in \mathcal{R}$. The optimization problem $\mathcal{P}$ is formulated as follows:
\begin{subequations} \label{eq:problem_formulation}
\begin{align}
\mathcal{P}: \max_{\{x_{ij}^{sd}\}} & \sum_{(s,d) \in \mathcal{R}} \frac{\sum_{(i,j)\in\mathcal{E}} x_{ij}^{sd} U_{ij}}{\sum_{(i,j)\in\mathcal{E}} x_{ij}^{sd}} \label{eq:obj} \\
\text{s.t.} \quad & \sum_{j} x_{sj}^{sd} - \sum_{j} x_{js}^{sd} = 1, \quad \forall (s,d)\in\mathcal{R} \label{eq:c1} \\
& \sum_{j} x_{dj}^{sd} - \sum_{j} x_{jd}^{sd} = -1, \quad \forall (s,d)\in\mathcal{R} \label{eq:c2} \\
& \sum_{j} x_{ij}^{sd} - \sum_{j} x_{ji}^{sd} = 0, \quad \forall i \neq s,d, \forall (s,d)\in\mathcal{R}. \label{eq:c3}
\end{align}
\end{subequations}
Constraints~\eqref{eq:c1} and~\eqref{eq:c2} enforce the flow requirements at the source and destination satellites, respectively. Constraint~\eqref{eq:c3} is the flow conservation condition for intermediate satellites.
